\chapter{Twenty-Two}\label{chapter-twenty-two}

\section{RMG Storage Systems (Git, IPLD,
GATOS)}\label{rmg-storage-systems-git-ipld-gatos}

\begin{quote}
State is Just a Frozen Worldline
\end{quote}

Every era of computing had to answer the question:

Where does the machine put things?

Punch cards, tapes, rotating platters, SSDs, object stores, CRDTs,
IPFS\emd{}they all tried. But every era made the same category mistake:

They stored bytes, not universes.

Storage was treated as something external to computation, an inert
substrate where results were dumped after-the-fact.

But in an RMG ontology, there is no difference between:

\begin{itemize}
\tightlist
\item
  state
\item
  history
\item
  computation
\item
  program
\end{itemize}

They are all the same thing: stable subgraphs of a rewrite universe.

This chapter unifies the world's most successful storage paradigm (Git),
the decentralized object graph protocol (IPLD), and the CΩMPUTER-native
design (GATOS) as three instantiations of the same idea:

Storage is computation, frozen. Computation is storage, evolving.

\begin{center}\rule{0.5\linewidth}{0.5pt}\end{center}

\subsection{1. State as Graphs, Not
Bytes}\label{state-as-graphs-not-bytes}

In classical systems, storage is trivially defined:

\begin{quote}
``Here are some bytes. Please don't lose them.''
\end{quote}

In CΩMPUTER, storage is the process of capturing a partial universe and
preserving its internal topology\emd{}its causal structure, its provenance,
its rewrite boundaries.

A ``file'' is:

\begin{itemize}
\tightlist
\item
  a sub-RMG
\item
  with canonicalized rewrite history
\item
  equipped with a consistent interface span
\item
  embedded in a larger cosmic graph
\end{itemize}

You don't store bytes. You store a stable region of the multiverse.

This turns storage into:

\begin{itemize}
\tightlist
\item
  a physics problem
\item
  a consistency problem
\item
  a rewrite problem
\end{itemize}

This is why Git, IPLD, and GATOS all converge in this chapter.

\begin{center}\rule{0.5\linewidth}{0.5pt}\end{center}

\subsection{2. Git as a Primitive RMG Storage
System}\label{git-as-a-primitive-rmg-storage-system}

Git accidentally stumbled into RMG theory 15 years before the math
existed.

Its core truths:

\begin{itemize}
\tightlist
\item
  Content-addressed objects → nodes in the RMG
\item
  Merkle edges → causal provenance
\item
  Commits → rewrite regions
\item
  Branches → worldlines
\item
  Merges → pushouts
\item
  Rebases → illicit time travel
\end{itemize}

Git is not a version control system.

Git is an early, naive, but shockingly successful RMG archiver.

Its flaws are exactly the places where it tries\emd{}and fails\emd{}to hide
its true nature:

\begin{itemize}
\tightlist
\item
  no real support for multi-parent universes
\item
  no explicit DPO/formal rewrite model
\item
  no structured semantics for conflicts
\item
  no native representation of rule sets
\item
  merges treated as ad hoc text manipulation
\end{itemize}

But the skeleton is unmistakable. Git is the ``Stone Age'' of RMG
storage\emd{}and it proved the model works.

\begin{center}\rule{0.5\linewidth}{0.5pt}\end{center}

\subsection{3. IPLD: The First Attempt at Universalizing the RMG
Layer}\label{ipld-the-first-attempt-at-universalizing-the-rmg-layer}

The InterPlanetary Linked Data (IPLD) effort recognized that Git's data
model was fundamental, but Git's text-first worldview was limiting.

IPLD generalized:

\begin{itemize}
\tightlist
\item
  arbitrary graphs
\item
  arbitrary codecs
\item
  arbitrary DAG semantics
\item
  stable content-addressing
\item
  decentralized storage
\end{itemize}

This was the first move toward a universal address space for RMG
fragments, but IPLD could not take the next step:

\begin{itemize}
\tightlist
\item
  no DPO rewrite semantics
\item
  no curvature tracking
\item
  no locality constraints
\item
  no dynamic rule sets
\item
  no multiworld execution model
\end{itemize}

IPLD provides the skeleton of a universal object graph. GATOS is the
musculature, ligaments, metabolism, and physics.

\begin{center}\rule{0.5\linewidth}{0.5pt}\end{center}

\subsection{4. GATOS: Storage as an Executable
Multiverse}\label{gatos-storage-as-an-executable-multiverse}

GATOS (Graph-As-The-Operating-Surface) is CΩMPUTER's native storage
substrate.

Its design principle:

\begin{quote}
Every stored object is a pocket universe with rules, invariants, and
possible futures.
\end{quote}

The filesystem is not a filesystem. It is a map of possible realities.

Each object has:

\begin{itemize}
\tightlist
\item
  Graph structure (the universe state)
\item
  Rewrite rules (the local physics)
\item
  Bundle indices (superposition structures)
\item
  Provenance paths (worldline history)
\item
  Constraints (laws)
\item
  Interfaces (spans into other universes)
\item
  Local curvature (optimization hints)
\end{itemize}

And crucially:

A GATOS object is executable.

Load it into the runtime, and the universe inside it evolves according
to its rule system.

This is the union of:

\begin{itemize}
\tightlist
\item
  Git's object model
\item
  IPLD's universality
\item
  CΩMPUTER's physics
\end{itemize}

GATOS is not a key-value store. It is a many-world hyper-database.

\begin{center}\rule{0.5\linewidth}{0.5pt}\end{center}

\subsection{5. The Core Formal Model: Objects as RMG
Patches}\label{the-core-formal-model-objects-as-rmg-patches}

Let's formalize.

A GATOS object O is defined as:

\begin{lstlisting}
O = (G,\, \mathcal{R},\, \mathcal{C},\, \mathcal{I},\, P,\, \mu)
\end{lstlisting}

Where:

\begin{itemize}
\tightlist
\item
  G --- Graph
\item
  $\mathcal{R}$ --- Rewrite rules
\item
  $\mathcal{C}$ --- Constraints
\item
  $\mathcal{I}$ --- Interface spans
\item
  P --- Provenance structure (worldline history)
\item
  $\mu$ --- Curvature metadata for optimization
\end{itemize}

Objects compose via span-based gluing:

\[
  O_1 \cup_{S} O_2 = \operatorname{Pushout}(O_1 \leftarrow S \rightarrow O_2)
\]

This is the categorical heart of the storage layer:

\begin{itemize}
\tightlist
\item
  merging is geometry
\item
  conflicts are curvature misalignments
\item
  rebases are illegal worldline rewrites
\item
  forks are multi-world splittings
\item
  snapshots are projections of the manifold
\end{itemize}

This isn't a metaphor: GATOS literally runs on these equations.

\begin{center}\rule{0.5\linewidth}{0.5pt}\end{center}

\begin{enumerate}
\def\labelenumi{\arabic{enumi}.}
\setcounter{enumi}{5}
\tightlist
\item
  Storage as Multiverse Folding: The Diagram
\end{enumerate}

Here's the core structural diagram you'll want in the book later:

\begin{lstlisting}
                      ┌──────────────────────┐
                      │        MRMW          │
                      │ (All Possible Worlds)│
                      └──────────┬───────────┘
                                 │
                  ┌──────────────┴──────────────┐
                  │            RMG               │
                  │    (Program Universe)        │
                  └──────────────┬──────────────┘
                                 │
                 ┌───────────────┴────────────────┐
                 │     Reachable Worldlines       │
                 └───────────────┬────────────────┘
                                 │
                  ┌──────────────┴──────────────┐
                  │      Stable Subgraphs        │
                  │          (State)             │
                  └──────────────┬──────────────┘
                                 │
         ┌───────────────────────┴─────────────────────────┐
         │            Storage Instantiations               │
         │                                                 │
┌────────┴───────────┐        ┌────────┴───────────┐       ┌────────┴──────────┐
│        Git         │        │        IPLD         │       │      GATOS        │
│   (Naive RMG)      │        │ (Generalized DAG)   │       │(Executable RMG OS)│
└────────────────────┘        └─────────────────────┘       └───────────────────┘
\end{lstlisting}

This is the chapter's money shot:

\begin{quote}
state emerges as stable subgraphs of the program universe, and
Git/IPLD/GATOS are simply different ways of freezing, naming, and
transmitting those subgraphs.
\end{quote}

\begin{center}\rule{0.5\linewidth}{0.5pt}\end{center}

\subsection{7. The Law of RMG Storage}\label{the-law-of-rmg-storage}

The first law of GATOS:

\begin{quote}
A stored object must preserve both structure (the graph) and possibility
(the rewrites).
\end{quote}

This means:

\begin{itemize}
\tightlist
\item
  a file is not merely data
\item
  a commit is not merely state
\item
  a checkpoint is not merely a snapshot
\end{itemize}

Everything stored in GATOS is:

\begin{enumerate}
\def\labelenumi{\arabic{enumi}.}
\tightlist
\item
  A worldline (how we got here)
\item
  A universe (where we are)
\item
  A rewrite horizon (where we can go from here)
\end{enumerate}

This is what it means to store computation-as-physics.

\begin{center}\rule{0.5\linewidth}{0.5pt}\end{center}

\subsection{8. Why This Matters for
CΩMPUTER}\label{why-this-matters-for-cux3c9mputer}

Without RMG storage systems:

\begin{itemize}
\tightlist
\item
  you cannot preserve provenance
\item
  you cannot re-enter a past universe
\item
  you cannot analyze curvature
\item
  you cannot do time-travel debugging
\item
  you cannot freeze or resume computation
\item
  you cannot run programs across multiple worlds
\item
  you cannot reason about counterfactuals
\item
  you cannot prove correctness across worldlines
\end{itemize}

Storage is the soul of CΩMPUTER. GATOS is the furnace that keeps that
soul alive.
